\lecture{6.1}{Software Setup and Simulated Data}
\label{lec:6-1}

\subsection{Required R Packages}

This chapter is a computational lab for the ideas developed in Chapter 5. The goal is
to see how cross-validated risk minimization and Super Learner ensembles are
implemented in practice. The examples use simulated data so that the true treatment
mechanism and true outcome regression are known.

The main package is \texttt{SuperLearner}. Several other packages are often useful
because they supply candidate algorithms or optimization routines. A typical setup is:

\begin{implementationbox}
\begin{lstlisting}[language=R]
pkgs <- c("SuperLearner", "quadprog", "nloptr", "MASS",
          "ggplot2", "nnls")

installed_packages <- row.names(installed.packages())

for (p in pkgs) {
  if (!(p %in% installed_packages)) {
    install.packages(p)
  }
  library(p, character.only = TRUE)
}
\end{lstlisting}
\end{implementationbox}

The packages \texttt{quadprog}, \texttt{nloptr}, and \texttt{nnls} are optimization
tools used by some Super Learner metalearners. The package \texttt{MASS} supplies
regression utilities, and \texttt{ggplot2} is useful for plotting. Additional candidate
algorithms may require additional packages.

\subsection{Loading \texttt{SuperLearner}}

After installation, the package can be loaded with:

\begin{implementationbox}
\begin{lstlisting}[language=R]
library(SuperLearner)
\end{lstlisting}
\end{implementationbox}

The package provides two kinds of functions:
\begin{itemize}
\item user-facing functions such as \texttt{SuperLearner()} and
\texttt{CV.SuperLearner()};
\item wrapper functions, for example \texttt{SL.glm} and \texttt{SL.mean};
these tell the package how to fit candidate algorithms.
\end{itemize}

To view available built-in learners and screeners:

\begin{implementationbox}
\begin{lstlisting}[language=R]
listWrappers(what = "SL")
listWrappers(what = "screen")
\end{lstlisting}
\end{implementationbox}

\subsection{Simulated Covariates}

We simulate $n$ independent observations with two covariates:
\[
W_1\sim N(0,1), \qquad W_2\sim \operatorname{Uniform}(0,1).
\]
These covariates are deliberately simple, but they are rich enough to illustrate
interactions and nonlinear candidate regressions.

\begin{implementationbox}
\begin{lstlisting}[language=R]
set.seed(212)

n <- 500
W1 <- rnorm(n = n)
W2 <- runif(n = n)
\end{lstlisting}
\end{implementationbox}

\subsection{Simulated Treatment Mechanism}

Treatment is binary. The true propensity score is
\[
g_1(W)=\Prob(A=1\given W)
=
\expit(-1+W_1-W_2+2W_1W_2).
\]
This treatment mechanism includes an interaction between $W_1$ and $W_2$. A main
terms logistic regression for $A$ on $W_1$ and $W_2$ is therefore misspecified.

\begin{implementationbox}
\begin{lstlisting}[language=R]
g1W <- plogis(-1 + W1 - W2 + 2 * W1 * W2)
A <- rbinom(n = n, size = 1, prob = g1W)
\end{lstlisting}
\end{implementationbox}

\begin{figure}[H]
\centering
\begin{tikzpicture}[
  node distance=1.5cm,
  var/.style={draw, circle, minimum size=.8cm},
  arr/.style={-{Latex[length=2.5mm]}, thick}
]
\node[var] (W1) {$W_1$};
\node[var, below=of W1] (W2) {$W_2$};
\node[var, right=2.2cm of W1] (A) {$A$};
\node[var, right=2.2cm of A] (Y) {$Y$};
\draw[arr] (W1) -- (A);
\draw[arr] (W2) -- (A);
\draw[arr] (W1) to[bend left=12] (Y);
\draw[arr] (W2) to[bend right=12] (Y);
\draw[arr] (A) -- (Y);
\node[draw, rounded corners=2pt, below=.55cm of A] (int) {$W_1W_2$};
\draw[arr] (int) -- (A);
\end{tikzpicture}
\caption{The simulated data have covariate-dependent treatment and outcome mechanisms.}
\label{fig:ch6-sim-dag}
\end{figure}

\subsection{Simulated Outcome Regression}

The true outcome regression is linear:
\[
\Qbar(A,W)=\E(Y\given A,W)=W_1-W_2+A.
\]
The observed outcome is generated by adding independent normal noise:
\[
Y=\Qbar(A,W)+\varepsilon,\qquad \varepsilon\sim N(0,1).
\]

\begin{implementationbox}
\begin{lstlisting}[language=R]
QAW <- W1 - W2 + A
Y <- QAW + rnorm(n = n)

full_data <- data.frame(W1 = W1, W2 = W2, A = A, Y = Y)
\end{lstlisting}
\end{implementationbox}

Because the true $\Qbar$ is known, we can check whether flexible methods behave
sensibly. However, in real applications the true regression function is unknown. The
point of the lab is to practice methods that are useful precisely because the truth is
not known.

\subsection{Reproducibility with Seeds}

Simulation, cross-validation splits, and some machine learning algorithms involve
randomness. A random seed makes the analysis reproducible:

\begin{implementationbox}
\begin{lstlisting}[language=R]
set.seed(212)
\end{lstlisting}
\end{implementationbox}

Using a seed does not make the result less random in a statistical sense. It fixes the
particular random draw so that another analyst running the same code obtains the
same simulated data, folds, and fitted results.

\begin{keyidea}
In a lab setting, seeds make examples reproducible. In a scientific analysis, they also
make computational results auditable: the same data and same code should lead to the
same reported estimates.
\end{keyidea}
